\subsection{Trace-preserving completions and physical channels}

\begin{theorem}[Trace-preserving completion outside a projection]
    \label{thm:matrix_support_completion_cptp}
    \lean{Matrix.supportCompletion_isKrausCPTP}
    \lean{Matrix.tracePrepareMap}
    \lean{Matrix.tracePrepareMap_apply}
    \lean{Matrix.tracePrepareMap_trace}
    \lean{Matrix.tracePrepareMap_isKrausCP}
    \lean{Matrix.tracePrepareMap_isKrausCPTP}
    \lean{Matrix.supportComplementMap}
    \lean{Matrix.supportComplementMap_apply}
    \lean{Matrix.supportComplementMap_isKrausCP}
    \lean{Matrix.supportComplementMap_trace}
    \lean{Matrix.supportCompletion}
    \lean{Matrix.supportCompletion_apply}
    \lean{Matrix.supportCompletion_isKrausCP}
    \lean{Matrix.supportCompletion_trace}
    \lean{Matrix.supportCompletion_apply_of_supported}
    \leanok
    \uses{def:is_kraus_cp, def:is_kraus_cptp,
        def:mpdo_state_preparation_map,
        def:mpdo_equiv_reindex_map,
        def:mpdo_right_partial_trace_map,
        def:mpdo_single_kraus_map}
    Let $P\in\MN{m}$ be an orthogonal projection, let
    $\rho\in\MN{n}$ be positive semidefinite with $\tr(\rho)=1$, and let
    $\mathcal E\colon\MN{m}\to\MN{n}$ be completely positive.  Define the
    trace-and-prepare map and its complementary restriction by
    \begin{align}
        \mathcal D_\rho(X)
        &=\tr(X)\rho,
        &
        \mathcal C_{P,\rho}(X)
        &=\mathcal D_\rho\bigl((I-P)X(I-P)^\dagger\bigr).
        \label{eq:matrix_support_complement_map}
    \end{align}
    Both maps are completely positive, and $\mathcal D_\rho$ is
    trace-preserving.  If
    \begin{align}
        \tr(\mathcal E(X))=\tr(PX)
        \label{eq:matrix_support_map_trace}
    \end{align}
    for every $X\in\MN{m}$, then the completion
    \begin{align}
        \widehat{\mathcal E}(X)
        &=\mathcal E(X)+\mathcal C_{P,\rho}(X)
        \label{eq:matrix_support_completion}
    \end{align}
    is trace-preserving and completely positive.  It agrees with
    $\mathcal E$ on every matrix supported by $P$:
    \begin{align}
        PXP=X
        \quad\Longrightarrow\quad
        \widehat{\mathcal E}(X)=\mathcal E(X).
        \label{eq:matrix_support_completion_agreement}
    \end{align}
\end{theorem}

\begin{proof}\leanok
    \uses{lem:mpdo_state_preparation_cp,
        lem:mpdo_state_preparation_cptp,
        thm:mpdo_equiv_reindex_cptp,
        lem:mpdo_right_partial_trace_cptp,
        lem:mpdo_single_kraus_map_apply,
        lem:mpdo_single_kraus_map_cp,
        lem:is_kraus_cp_comp,
        lem:is_kraus_cptp_comp,
        lem:is_kraus_cp_add,
        thm:is_kraus_cptp_of_is_kraus_cp_trace_preserving}
    The map $\mathcal D_\rho$ factors as state preparation, interchange of
    the two tensor factors, and partial trace.  The map
    $X\mapsto(I-P)X(I-P)^\dagger$ has a single Kraus operator.  Thus the two
    maps in~(\ref{eq:matrix_support_complement_map}) are completely positive.
    Since $P=P^\dagger=P^2$,
    \begin{align}
        \tr\bigl(\mathcal C_{P,\rho}(X)\bigr)
        &=\tr((I-P)X).
        \notag
    \end{align}
    Adding this identity to~(\ref{eq:matrix_support_map_trace}) gives
    $\tr(\widehat{\mathcal E}(X))=\tr(X)$.  Finally, $PXP=X$ implies
    $(I-P)X(I-P)=0$, which proves
    (\ref{eq:matrix_support_completion_agreement}).
\end{proof}

\begin{theorem}[Physical channels from unitary sector conjugacies]
    \label{thm:mpdo_bnt_conditional_physical_channels}
    \lean{MPOTensor.BNTAlgebraTensorClause.TwoSiteExactSectorGauge.%
        UnitarySectorConjugacy.isRFPViaTS}
    \lean{MPOTensor.BNTAlgebraTensorClause.oneSiteRetainedProjection}
    \lean{MPOTensor.BNTAlgebraTensorClause.oneSiteRetainedProjection_isHermitian}
    \lean{MPOTensor.BNTAlgebraTensorClause.oneSiteRetainedProjection_mul_self}
    \lean{MPOTensor.BNTAlgebraTensorClause.TwoSiteMultiplicitySpectrum.%
        twoSiteRetainedProjection}
    \lean{MPOTensor.BNTAlgebraTensorClause.TwoSiteMultiplicitySpectrum.%
        twoSiteRetainedProjection_isHermitian}
    \lean{MPOTensor.BNTAlgebraTensorClause.TwoSiteMultiplicitySpectrum.%
        twoSiteRetainedProjection_mul_self}
    \lean{MPOTensor.BNTAlgebraTensorClause.TwoSiteExactSectorGauge.%
        UnitarySectorConjugacy.rawPhysicalT}
    \lean{MPOTensor.BNTAlgebraTensorClause.TwoSiteExactSectorGauge.%
        UnitarySectorConjugacy.rawPhysicalS}
    \lean{MPOTensor.BNTAlgebraTensorClause.TwoSiteExactSectorGauge.%
        UnitarySectorConjugacy.rawPhysicalT_isKrausCP}
    \lean{MPOTensor.BNTAlgebraTensorClause.TwoSiteExactSectorGauge.%
        UnitarySectorConjugacy.rawPhysicalS_isKrausCP}
    \lean{MPOTensor.BNTAlgebraTensorClause.TwoSiteExactSectorGauge.%
        UnitarySectorConjugacy.trace_rawPhysicalT}
    \lean{MPOTensor.BNTAlgebraTensorClause.TwoSiteExactSectorGauge.%
        UnitarySectorConjugacy.trace_rawPhysicalS}
    \lean{MPOTensor.BNTAlgebraTensorClause.TwoSiteExactSectorGauge.%
        UnitarySectorConjugacy.rawPhysicalT_physClose1}
    \lean{MPOTensor.BNTAlgebraTensorClause.TwoSiteExactSectorGauge.%
        UnitarySectorConjugacy.rawPhysicalS_physClose2}
    \lean{MPOTensor.BNTAlgebraTensorClause.TwoSiteExactSectorGauge.%
        oneSiteRetainedProjection_physClose1}
    \lean{MPOTensor.BNTAlgebraTensorClause.TwoSiteExactSectorGauge.%
        twoSiteRetainedProjection_physClose2}
    \lean{MPOTensor.BNTAlgebraTensorClause.TwoSiteExactSectorGauge.%
        UnitarySectorConjugacy.physicalT}
    \lean{MPOTensor.BNTAlgebraTensorClause.TwoSiteExactSectorGauge.%
        UnitarySectorConjugacy.physicalS}
    \lean{MPOTensor.BNTAlgebraTensorClause.TwoSiteExactSectorGauge.%
        UnitarySectorConjugacy.physicalT_isKrausCPTP}
    \lean{MPOTensor.BNTAlgebraTensorClause.TwoSiteExactSectorGauge.%
        UnitarySectorConjugacy.physicalS_isKrausCPTP}
    \lean{MPOTensor.BNTAlgebraTensorClause.TwoSiteExactSectorGauge.%
        UnitarySectorConjugacy.physicalT_physClose1}
    \lean{MPOTensor.BNTAlgebraTensorClause.TwoSiteExactSectorGauge.%
        UnitarySectorConjugacy.physicalS_physClose2}
    \leanok
    \uses{def:mpdo_bnt_ambient_sector_maps,
        def:mpdo_bnt_direct_sum_unitary_conjugations,
        def:mpdo_faithful_uniform_density,
        def:is_rfp_via_ts}
    Suppose that the one-site and two-site sector decompositions are paired by
    a relabelling $\sigma$ and unitaries $U_\gamma$ satisfying
    \begin{align}
        A_{\sigma(\gamma)}^i
        &=U_\gamma M_\gamma^i U_\gamma^\dagger
        \label{eq:mpdo_bnt_physical_channel_unitary_conjugacy}
    \end{align}
    for every sector $\gamma$ and every vertical letter $i$.  With the maps of
    the preceding two subsections, set
    \begin{align}
        \mathcal T_0
        &=\widetilde R_2\circ\widetilde T\circ R_1,
        &
        \mathcal S_0
        &=\widetilde R_1\circ\widetilde S\circ R_2.
        \label{eq:mpdo_bnt_active_physical_maps}
    \end{align}
    Let
    \begin{align}
        P_1&=W_1^\dagger W_1,
        &
        P_2&=\kappa_\ast(W_2^\dagger W_2),
        &
        Q_i&=I-P_i.
        \label{eq:mpdo_bnt_physical_active_projections}
    \end{align}
    The matrices $P_1$ and $P_2$ are orthogonal projections.  If $d>0$, put
    $\tau_1=d^{-1}I_d$ and $\tau_2=d^{-2}I_{d^2}$, and define
    \begin{align}
        \mathcal T(X)
        &=\mathcal T_0(X)+\tr(Q_1XQ_1)\tau_2,
        \notag\\
        \mathcal S(Y)
        &=\mathcal S_0(Y)+\tr(Q_2YQ_2)\tau_1.
        \label{eq:mpdo_bnt_completed_physical_maps}
    \end{align}
    If $d=0$, take both maps to be zero; their domain and codomain matrix
    algebras then consist only of zero.
    Then $\mathcal T$ and $\mathcal S$ are trace-preserving completely positive
    maps on the full physical matrix algebras.  For every virtual matrix $Z$,
    they satisfy
    \begin{align}
        \mathcal T[M_1(Z)]&=M_2(Z),
        &
        \mathcal S[M_2(Z)]&=M_1(Z).
        \label{eq:mpdo_bnt_completed_physical_map_actions}
    \end{align}
    Hence $M$ is a renormalization fixed point in the sense of
    Definition~\ref{def:is_rfp_via_ts}.

    The compositions in
    (\ref{eq:mpdo_bnt_active_physical_maps}) are the formulas of
    \cite[Appendix~C.4, lines~2065--2085]{Cirac2016MPDO_arXiv}.  Those formulas
    determine the action after compression to the retained physical sectors.
    With rectangular coisometries, the literal composites annihilate the
    discarded zero-sector complements and fail to preserve trace there.
    Definition~4.1 nevertheless requires maps on the full physical algebras
    \cite[Definition~4.1, lines~638--660]{Cirac2016MPDO_arXiv}.
    The source does not specify a trace-restoring modification on the discarded
    complements.  Formula~(\ref{eq:mpdo_bnt_completed_physical_maps}) makes one
    such choice when the physical space is nonzero.  Replacing
    $\tau_1,\tau_2$ by any density matrices gives the same action on the
    physical operators generated by $M$.  The conclusion is conditional on
    the unitaries in
    (\ref{eq:mpdo_bnt_physical_channel_unitary_conjugacy}); their existence is
    not asserted here.
\end{theorem}

\begin{proof}\leanok
    \uses{thm:mpdo_bnt_ambient_sector_complete_positivity,
        thm:mpdo_bnt_ambient_sector_trace,
        thm:mpdo_bnt_ambient_sector_tensor_letters,
        thm:mpdo_bnt_direct_sum_unitary_conjugations,
        thm:matrix_support_completion_cptp,
        thm:mpdo_faithful_uniform_density,
        lem:mpdo_coisometric_compression_trace,
        lem:mpdo_vertical_bond_contraction_phys_close,
        lem:mpdo_vertical_bond_contraction_reconstruction,
        thm:mpdo_one_site_closure_two_site_blocking}
    Complete positivity of the four outer maps and of the two direct-sum
    conjugations gives complete positivity of $\mathcal T_0$ and
    $\mathcal S_0$.  Their traces are
    \begin{align}
        \tr(\mathcal T_0(X))&=\tr(P_1X),
        &
        \tr(\mathcal S_0(Y))&=\tr(P_2Y).
        \label{eq:mpdo_bnt_active_physical_map_traces}
    \end{align}
    Suppose first that $d>0$.  Since $Q_i$ is a projection and $\tau_i$ has
    trace one, the corresponding complementary term is completely positive
    and has trace $\tr(Q_iXQ_i)=\tr(Q_iX)$.  Adding
    (\ref{eq:mpdo_bnt_active_physical_map_traces}) and using $P_i+Q_i=I$
    proves trace preservation of the completed maps.  If $d=0$, every ambient
    matrix is zero, so the zero maps are completely positive and trace
    preserving.

    The reconstruction identities for the two vertical decompositions imply
    \begin{align}
        P_1M_1(Z)P_1&=M_1(Z),
        &
        P_2M_2(Z)P_2&=M_2(Z).
        \label{eq:mpdo_bnt_physical_closure_support}
    \end{align}
    Thus both complementary terms vanish on the corresponding physical
    operators.  The retraction formulas extract the families
    $(m_\gamma M_\gamma^i)_\gamma$ and
    $(m_\gamma A_{\sigma(\gamma)}^i)_\gamma$; unitary conjugation carries
    these families into one another, and the normalized restorations recover
    the one-site and two-site letters.  Linearity in the virtual matrix gives
    (\ref{eq:mpdo_bnt_completed_physical_map_actions}).
\end{proof}

\begin{theorem}[Renormalization fixed point from positive-tail reflected
    sector targets]
    \label{thm:mpdo_bnt_conditional_rfp_from_reflected_targets}
    \lean{MPOTensor.BNTAlgebraTensorClause.TwoSiteExactSectorGauge.%
        isRFPViaTS_of_positive_tail_reflected_target}
    \lean{MPOTensor.BNTAlgebraTensorClause.TwoSiteExactSectorGauge.%
        UnitarySectorConjugacy.ofPositiveTailReflectedTarget}
    \leanok
    \uses{def:mpdo, def:cpsv_canonical_form,
        def:mpdo_bnt_two_site_exact_sector_gauge,
        def:is_rfp_via_ts}
    Let the one-site and blocked two-site BNT decompositions have the
    source-derived exact sector pairing, including the matched multiplicity
    spectrum, and suppose that $M$ generates matrix product density operators
    and that its doubled-index tensor is in literal CPSV canonical form.
    Assume, for every sector $\gamma$, every positive tail length, every pair
    of open bond indices, and every pair of two-site tail words, that the
    identity-dressed marked-chain coefficient equals the corresponding
    reflected-adjoint marked-chain coefficient of $M_\gamma$ in the adjoint
    two-site blocking, with both tail words reversed on the reflected-adjoint
    side. This is the additional comparison needed to reproduce the Figure~8
    argument in
    \cite[Proposition~4.13, lines~1903--1921]{Cirac2016MPDO_arXiv} at its
    invocation in Appendix~C.4, line~2057. Then there are trace-preserving
    completely positive maps $\mathcal T$ and $\mathcal S$ satisfying the
    identities in
    (\ref{eq:mpdo_bnt_completed_physical_map_actions}) for every virtual
    matrix $X$. Hence $M$ is a renormalization fixed point in the sense of
    Definition~\ref{def:is_rfp_via_ts}.
    The target is an explicit hypothesis of this theorem.  It is derived from
    the BNT algebra clause by the mixed-prefix comparison in
    Theorem~\ref{thm:mpdo_bnt_mixed_prefix_unitary_sector_conjugacy}.
\end{theorem}

\begin{proof}\leanok
    \uses{thm:mpdo_bnt_conditional_identity_marked_unitary_conjugacy,
        thm:mpdo_bnt_conditional_physical_channels}
    For each sector,
    Theorem~\ref{thm:mpdo_bnt_conditional_identity_marked_unitary_conjugacy}
    supplies a unitary $U_\gamma$ satisfying
    (\ref{eq:mpdo_bnt_conditional_unitary_conjugacy}), which is precisely
    (\ref{eq:mpdo_bnt_physical_channel_unitary_conjugacy}). Choose these
    unitaries simultaneously over the finite set of sectors.
    Theorem~\ref{thm:mpdo_bnt_conditional_physical_channels} then gives
    $\mathcal T$ and $\mathcal S$ together with
    (\ref{eq:mpdo_bnt_completed_physical_map_actions}).
\end{proof}

\begin{theorem}[BNT algebra clause implies a renormalization fixed point]
    \label{thm:mpdo_bnt_algebra_tensor_to_rfp}
    \lean{MPOTensor.BNTAlgebraTensorClause.isRFPViaTS}
    \lean{MPOTensor.HasBNTAlgebraTensorClause.isRFPViaTS}
    \leanok
    \uses{def:mpdo, def:cpsv_canonical_form,
        def:mpdo_bnt_algebra_tensor_clause, def:is_rfp_via_ts}
    Let $M$ be an MPDO whose doubled-index tensor is in literal CPSV
    canonical form.  If $M$ has a tensor-attached BNT algebra clause, then
    there are trace-preserving completely positive maps
    \begin{align}
        \mathcal T\colon\MN d&\longrightarrow\MN{d^2},
        &
        \mathcal S\colon\MN{d^2}&\longrightarrow\MN d
        \notag
    \end{align}
    such that, for every virtual matrix $X$,
    \begin{align}
        \mathcal T[M_1(X)]&=M_2(X),
        &
        \mathcal S[M_2(X)]&=M_1(X).
        \label{eq:mpdo_bnt_algebra_tensor_rfp_maps}
    \end{align}
    Thus $M$ is a renormalization fixed point in the sense of
    Definition~\ref{def:is_rfp_via_ts}.  This is implication
    (ii)$\Rightarrow$(i) of
    \cite[Theorem~4.14, lines~972--993]{Cirac2016MPDO_arXiv}.

    \textbf{Local fix (mixed one-site/two-site prefix):} The common marked
    comparison needed to normalize the exact sector gauges is supplied by
    Theorem~\ref{thm:mpdo_bnt_mixed_prefix_unitary_sector_conjugacy}, as
    recorded in
    \path{docs/paper-gaps/cpsv16_two_site_sector_unitary_gauge_gap.tex}.

    \textbf{Local fix (zero-sector complement):} The physical maps are
    completed on the discarded complements as recorded in
    \path{docs/paper-gaps/cpgsv17_vertical_isometry_zero_sector.tex}.
\end{theorem}

\begin{proof}\leanok
    \uses{thm:mpdo_bnt_algebra_tensor_clause_spectrum,
        thm:mpdo_bnt_mixed_prefix_unitary_sector_conjugacy,
        thm:mpdo_bnt_conditional_physical_channels}
    The tensor-attached algebra clause determines the two-site vertical
    decomposition, the sector relabelling, and exact invertible gauges.
    Theorem~\ref{thm:mpdo_bnt_mixed_prefix_unitary_sector_conjugacy}
    normalizes these gauges simultaneously to unitaries.  Apply
    Theorem~\ref{thm:mpdo_bnt_conditional_physical_channels} to obtain the two
    trace-preserving completely positive maps and
    (\ref{eq:mpdo_bnt_algebra_tensor_rfp_maps}).
\end{proof}
